\section{Apprentissages scientifiques et techniques}

\begin{guideline}
C'est une section centrale du rapport SRS. Elle doit démontrer votre
capacité à acquérir et maîtriser des concepts techniques et scientifiques
dans un contexte professionnel réel. Pour chaque apprentissage
significatif, expliquez le contexte qui l'a rendu nécessaire et ce que
vous en retenez. Soyez précis et technique : c'est ici que vous prouvez
votre niveau d'ingénieur.
\end{guideline}


\subsection{Concepts techniques et scientifiques acquis}

\begin{guideline}
Détaillez les algorithmes, protocoles, architectures, frameworks ou
standards que vous avez dû apprendre ou approfondir. Pour chaque concept,
expliquez pourquoi il était nécessaire dans votre contexte de stage.
\end{guideline}

\instruction{[Pour chaque concept technique majeur, rédigez un
paragraphe structuré ainsi : (1) quel était le besoin concret qui a
rendu cet apprentissage nécessaire, (2) ce que vous avez appris et
comment vous l'avez mis en œuvre, (3) ce que vous en retenez. Ex. :
protocoles OpenSSH sur Windows Server, architecture AWX/Ansible,
gestion des Execution Environments, sécurité de la chaîne
d'approvisionnement logicielle avec Nexus...]}


\subsection{Méthodes et processus mobilisés}

\begin{guideline}
Décrivez les méthodologies de travail, démarches qualité ou de gestion
de projet que vous avez appliquées : méthode agile, CI/CD, GitOps,
revues de code, tests d'intégration, documentation technique...
\end{guideline}

\instruction{[Expliquez les méthodes et processus que vous avez
appliqués au quotidien. Distinguez ce que vous connaissiez déjà de ce
que vous avez découvert en stage. Ex. : approche GitOps pour la
gestion de la dérive de configuration, structuration des rôles Ansible,
utilisation de Docker et ansible-builder pour les Execution
Environments...]}


\subsection{Aspects juridiques et réglementaires}

\begin{guideline}
Mentionnez les contraintes légales ou réglementaires que vous avez
rencontrées : RGPD, licences open source, conformité sectorielle,
politiques de sécurité internes, normes ISO...
\end{guideline}

\instruction{[Décrivez les aspects juridiques ou réglementaires qui ont
influencé vos choix techniques. Ex. : gestion des licences des
collections Ansible Galaxy, politique de sécurité Azure, contraintes
de confidentialité sur les données de configuration, conformité aux
politiques de least privilege...]}